\section{De la síntesis inductiva a la IA del mundo físico}

Gao Jiyang (高继扬) vuelve una y otra vez a un mismo método: inducir y sintetizar. La primaria, la olimpiada de física, las publicaciones del doctorado, las entrevistas de trabajo y la decisión de emprender las aborda de la misma manera: «primero fijar el objetivo, luego descomponer el camino y después aumentar la probabilidad de acertar». El método parece sencillo, pero explica por qué su trayectoria lo llevó del Departamento de Ingeniería Electrónica de la Universidad Tsinghua (清华), pasando por el aprendizaje profundo, la visión por computadora y la conducción autónoma, hasta la inteligencia corpórea.

\begin{figure}[H]
\centering
\includegraphics[width=0.96\textwidth]{trajectory-industrial-training.png}
\caption{La trayectoria de formación industrial de Gao Jiyang: síntesis inductiva, iniciación en la IA, Waymo, Momenta y Galaxea (星海图).}
\end{figure}

\begin{knowledgebox}{Lectura de la figura: cómo se acumulan tres capacidades}
La ruta de la figura no es una lista de cargos. La olimpiada de física y el doctorado le enseñaron a inducir, descomponer y gestionar probabilidades; las prácticas en SenseTime (商汤) le permitieron sentir por primera vez que una red neuronal puede extraer regularidades de los datos por sí sola; Waymo le dio formación en ingeniería de sistemas y en medición; Momenta, en la entrega de productos y el valor para el cliente; y Galaxea lleva todas estas capacidades a la inteligencia corpórea.
\end{knowledgebox}

\subsection{Zeng Guofan: del letrado purista a los logros prácticos}

La sección anterior trató la «síntesis inductiva» de Gao Jiyang; esta explica de dónde viene su realismo. El ejemplo de Zeng Guofan (曾国藩) no es un adorno histórico, sino uno de los primeros modelos con los que entendió la organización de una empresa emergente.

En la entrevista, Gao Jiyang cuenta que leyó a Zeng Guofan después de un tropiezo al solicitar ingreso a las universidades. No le interesaba la curiosidad histórica, sino «cómo un letrado confuciano purista se convierte en alguien capaz de movilizar recursos, organizar equipos y librar batallas difíciles en el mundo real». Este pasaje es fundamental para entender su visión del emprendimiento: para hacer cosas en el mundo real no bastan las ideas, los artículos ni la inteligencia; hay que movilizar recursos, organizar personas y asumir los resultados.

\begin{importantbox}{La perspectiva de los logros prácticos}
Los logros prácticos no son utilitarismo de corto plazo, sino convertir un objetivo en recursos, organización, acciones y resultados. Un emprendimiento de robótica necesita especialmente esta perspectiva, porque cada uno de sus pasos tiene que materializarse en el mundo real.
\end{importantbox}

\subsection{El atractivo del aprendizaje profundo}

Durante sus prácticas en SenseTime entrenó una red neuronal por primera vez y percibió la fascinación de que «la máquina extraiga regularidades de los datos por sí misma». En la programación tradicional, una persona resume las reglas y luego las escribe como código; en el aprendizaje automático, las reglas se condensan en los parámetros y el modelo aprende a partir de la distribución de los datos. Ese momento fijó su manera de entender la IA en adelante: el valor de la IA no está solo en automatizar, sino en cambiar la forma en que los seres humanos descubren regularidades y construyen productividad.

\begin{knowledgebox}{Términos clave: de las reglas al aprendizaje}
\begin{tabularx}{\textwidth}{p{0.22\textwidth}p{0.32\textwidth}X}
\toprule
Término & Problema que resuelve & Significado en este episodio \\
\midrule
Rule-based & Reglas, módulos e if-else escritos a mano por personas & Base importante de los sistemas tradicionales de conducción autónoma y de robótica. \\
Data-driven & Aprender regularidades a partir de los datos & La razón central por la que Gao Jiyang eligió la IA y la conducción autónoma. \\
Benchmark & Medir el desempeño de los modelos con una evaluación unificada & Lo necesitan la publicación de artículos, la iteración de modelos y la recipe de datos. \\
Physical World AI & IA que actúa directamente sobre tareas del mundo físico & La conducción autónoma es su primera forma; la inteligencia corpórea, la siguiente. \\
\bottomrule
\end{tabularx}
\end{knowledgebox}

\subsection{Resumen de la sección}

El hilo conductor de Gao Jiyang es convertir el esfuerzo personal en un método reutilizable y luego llevar la capacidad de la IA basada en datos al mundo físico. Lo que le atrae de la conducción autónoma y de los robots es precisamente que la IA se vuelve una variable de base de toda la industria, no una herramienta de optimización local.
